\documentclass[11pt, letterpaper]{article}

% Fonts and Encoding
\usepackage[utf8]{inputenc}
\usepackage[T1]{fontenc}
\usepackage{lmodern} % High-quality Latin Modern vector font

% Math and Symbols
\usepackage{amsmath, amssymb, amsthm, amsfonts, mathrsfs, mathtools}

% Layout and Geometry
\usepackage{geometry}
\geometry{letterpaper, margin=1in}
\setlength{\parindent}{0pt}
\setlength{\parskip}{0.6em}

% Formatting and Graphics
\usepackage{graphicx, xcolor, enumitem, booktabs, array, tabularx}

% TikZ for Illustrations
\usepackage{tikz}
\usetikzlibrary{arrows.meta, positioning, shapes.geometric, calc, backgrounds, matrix}

% Clickable Links and TOC
\usepackage{hyperref}
\hypersetup{
    colorlinks=true,
    linkcolor=blue,
    filecolor=magenta,      
    urlcolor=cyan,
    pdftitle={Transforms of Periodic and Power Functions},
}

% Theorem Environments
\newtheorem{theorem}{Theorem}[section]
\newtheorem{lemma}[theorem]{Lemma}
\newtheorem{proposition}[theorem]{Proposition}
\newtheorem{corollary}[theorem]{Corollary}
\theoremstyle{definition}
\newtheorem{definition}{Definition}[section]
\newtheorem{axiom}{Axiom}[section]
\newtheorem{example}{Example}[section]
\theoremstyle{remark}
\newtheorem{remark}{Remark}[section]

% Custom Commands (Enforced standard notation)
\newcommand{\R}{\mathbb{R}}
\newcommand{\N}{\mathbb{N}}
\newcommand{\Z}{\mathbb{Z}}
\newcommand{\Q}{\mathbb{Q}}
\newcommand{\E}{\mathbb{E}}
\newcommand{\Prob}{\mathbb{P}}
\newcommand{\F}{\mathcal{F}}

\begin{document}

\title{\textbf{Transforms of Periodic and Power Functions\\ \Large Comprehensive Problem Sets and Solutions}}
\author{Tashfeen Omran}
\date{November 5, 2026}
\maketitle

\tableofcontents
\newpage

\section{Concept Checklist \& Question Mapping}

The following table maps the core concepts extracted from the source material to the corresponding problems in this document. This ensures comprehensive coverage of all techniques, algebraic manipulations, and theoretical applications.

\renewcommand{\arraystretch}{1.5}
\begin{tabularx}{\textwidth}{@{} l X >{$\displaystyle}c<{$} @{}}
\toprule
\textbf{ID} & \textbf{Concept / Problem Type} & \text{\textbf{Relevant Questions}} \\
\midrule
C1 & Direct Periodic Transforms (Square, Sawtooth, Triangular) & 1, 2, 3, 7, 8 \\
C2 & Rectified Trigonometric Waves (Half/Full Rectification) & 4, 5, 9, 10 \\
C3 & Piecewise Exponential and Polynomial Periodic Functions & 6, 11, 12 \\
C4 & Inverse Transforms via Geometric Series Expansions & 13, 14, 15, 16, 17 \\
C5 & Gamma Function Evaluation and Recursive Properties & 18, 19, 22 \\
C6 & Transforms of Fractional Power Functions ($t^r$) & 20, 21, 23 \\
C7 & Formal Series Expansions for Transforms (Maclaurin) & 24, 25, 26 \\
C8 & Initial Value Problems with Periodic Forcing Functions & 27, 28 \\
C9 & Dirac Comb (Impulse Train) and Derivative Shortcuts & 29 \\
C10 & Diagnostic Testing: "Find the Flaw" & 30 \\
\bottomrule
\end{tabularx}

\vspace{1em}
\begin{remark}
The problems are designed to progressively build intuition, starting from direct integration and culminating in abstract series expansions and differential equations.
\end{remark}

\newpage
\section{Problem Sets}

\subsection{Direct Periodic Transforms}
\textbf{1.} Determine $\mathcal{L}\{f\}$, where $f(t)$ is a square wave defined by $f(t) = 3$ for $0 < t < 2$, and $f(t) = 0$ for $2 < t < 4$, with period $T = 4$.

\textbf{2.} Determine $\mathcal{L}\{f\}$, where $f(t) = 3t$ for $0 < t < 3$, and $f(t)$ has period $T = 3$.

\textbf{3.} Determine $\mathcal{L}\{f\}$, where $f(t)$ is a triangular wave defined by $f(t) = 2t$ for $0 < t < 1$, and $f(t) = 4 - 2t$ for $1 < t < 2$, with period $T = 2$.

\textbf{4.} Determine $\mathcal{L}\{f\}$, where $f(t)$ is a half-rectified sine wave defined by $f(t) = \sin(3t)$ for $0 < t < \pi/3$, and $f(t) = 0$ for $\pi/3 < t < 2\pi/3$, with period $T = 2\pi/3$.

\textbf{5.} Determine $\mathcal{L}\{f\}$, where $f(t) = |\cos(t)|$. (Hint: Determine the fundamental period first).

\textbf{6.} Determine $\mathcal{L}\{f\}$, where $f(t) = e^{2t}$ for $0 < t < 1$, with period $T = 1$.

\textbf{7.} Determine $\mathcal{L}\{f\}$, where $f(t) = t^2$ for $0 < t < 2$, with period $T = 2$.

\textbf{8.} Determine $\mathcal{L}\{f\}$, where $f(t) = 1$ for $0 < t < a$, and $f(t) = -1$ for $a < t < 2a$, with period $T = 2a$.

\textbf{9.} Determine $\mathcal{L}\{f\}$, where $f(t) = \sin(t)$ for $0 < t < \pi$, and $f(t) = \sin(t)$ for $\pi < t < 2\pi$ is not true, rather $f(t) = |\sin(t)|$ (Full-rectified sine wave).

\textbf{10.} Determine $\mathcal{L}\{f\}$, where $f(t) = \sinh(t)$ for $0 < t < 2$, with period $T = 2$.

\textbf{11.} Determine $\mathcal{L}\{f\}$, where $f(t) = t e^{-t}$ for $0 < t < 1$, with period $T = 1$.

\textbf{12.} Determine $\mathcal{L}\{f\}$, where $f(t) = 1 - e^{-t}$ for $0 < t < 2$, with period $T = 2$.

\subsection{Inverse Transforms via Geometric Series}
\textbf{13.} Find the inverse Laplace transform of $F(s) = \frac{1}{s(1 - e^{-4s})}$. Express the result as an infinite series of step functions.

\textbf{14.} Find the inverse Laplace transform of $F(s) = \frac{1}{(s^2 + 9)(1 - e^{-\pi s})}$.

\textbf{15.} Find the inverse Laplace transform of $F(s) = \frac{e^{-2s}}{(s + 3)(1 - e^{-2s})}$.

\textbf{16.} Find the inverse Laplace transform of $F(s) = \frac{1 - e^{-s}}{s^2(1 - e^{-2s})}$.

\textbf{17.} Show that $\mathcal{L}^{-1} \left\{ \frac{1}{s(1 + e^{-s})} \right\}$ represents an alternating square wave.

\subsection{Gamma Function and Fractional Powers}
\textbf{18.} Evaluate $\Gamma(9/2)$ using the recursive property of the Gamma function.

\textbf{19.} Evaluate the integral $\int_0^\infty x^4 e^{-x^2} \,dx$ by relating it to the Gamma function.

\textbf{20.} Determine $\mathcal{L}\{t^{7/2}\}$.

\textbf{21.} Determine $\mathcal{L}\{t^{-1/2} e^{4t}\}$ using the First Translation Theorem.

\textbf{22.} Show that $\Gamma(n + 1/2) = \frac{(2n)! \sqrt{\pi}}{4^n n!}$ for any integer $n \ge 0$.

\textbf{23.} Find the inverse Laplace transform of $F(s) = s^{-9/2}$.

\subsection{Formal Series Expansions}
\textbf{24.} Find an expansion for $\mathcal{L} \left\{ \frac{\sinh(t)}{t} \right\}$ in powers of $1/s$ by transforming the Maclaurin series term-by-term.

\textbf{25.} Find an expansion for $\mathcal{L} \{\cos(2\sqrt{t})\}$ by transforming its Maclaurin series term-by-term.

\textbf{26.} Find the inverse Laplace transform of $F(s) = \frac{1}{s} e^{-3/s}$ using a series expansion.

\subsection{Differential Equations and Advanced Models}
\textbf{27.} Solve the initial value problem $y' + 2y = f(t)$ with $y(0) = 0$, where $f(t)$ is a square wave defined by $f(t) = 1$ for $0 < t < 1$ and $f(t) = 0$ for $1 < t < 2$, with period $T = 2$.

\textbf{28.} Solve the initial value problem $y'' + y = \text{III}_{2\pi}(t)$ with $y(0) = 0, y'(0) = 0$, where $\text{III}_{2\pi}(t) = \sum_{n=0}^\infty \delta(t - 2n\pi)$ is the Dirac comb.

\textbf{29.} Use the derivative shortcut $\mathcal{L}\{f'\} = sF(s) - f(0)$ to find the Laplace transform of the sawtooth wave $f(t) = t$ for $0 < t < a$, period $a$. (Hint: Express $f'(t)$ using the Dirac comb).

\subsection{Diagnostic Testing: "Find the Flaw"}
\textbf{30.} \textbf{Flawed Problem:} Find the Laplace transform of $f(t) = t^{3/2}$. \\
\textit{Flawed Solution:} Using the formula $\mathcal{L}\{t^r\} = \frac{\Gamma(r+1)}{s^{r+1}}$, we have $r = 3/2$.
\[ \mathcal{L}\{t^{3/2}\} = \frac{\Gamma(3/2)}{s^{3/2}} \]
We know $\Gamma(3/2) = \frac{1}{2}\Gamma(1/2) = \frac{1}{2}\sqrt{\pi}$. Therefore, the final answer is $\frac{\sqrt{\pi}}{2s^{3/2}}$. \\
\textbf{Task:} Identify the error, explain why it is wrong in one sentence, and provide the correct final answer.

\newpage
\section{Detailed Solutions}

\textbf{Problem 1:} Determine $\mathcal{L}\{f\}$, where $f(t)$ is a square wave defined by $f(t) = 3$ for $0 < t < 2$, and $f(t) = 0$ for $2 < t < 4$, with period $T = 4$.

\textbf{Solution:}
The period is $T = 4$. The formula for the Laplace transform of a periodic function is:
\[ \mathcal{L}\{f\}(s) = \frac{1}{1 - e^{-4s}} \int_0^4 e^{-st} f(t) \,dt \]
Since $f(t) = 0$ on $(2, 4)$, we restrict the integral:
\[ \int_0^2 3 e^{-st} \,dt = 3 \left[ -\frac{1}{s} e^{-st} \right]_0^2 = \frac{3}{s} (1 - e^{-2s}) \]
Substitute this back into the periodic transform formula:
\[ \mathcal{L}\{f\}(s) = \frac{3(1 - e^{-2s})}{s(1 - e^{-4s})} \]
Notice that $1 - e^{-4s}$ is a difference of squares: $(1 - e^{-2s})(1 + e^{-2s})$. Canceling the common term yields:
\[ \mathcal{L}\{f\}(s) = \frac{3}{s(1 + e^{-2s})} \]

\begin{center}
\begin{tikzpicture}[scale=1.2]
    \fill[white] (-1,-1) rectangle (7,4); % Background layer equivalent
    \draw[thick, ->, >=stealth] (-0.5, 0) -- (6, 0) node[right] {$t$};
    \draw[thick, ->, >=stealth] (0, -0.5) -- (0, 3.5) node[above] {$f(t)$};
    
    \draw[ultra thick, blue!80!black] (0, 3) -- (2, 3);
    \draw[ultra thick, blue!80!black, dashed] (2, 3) -- (2, 0);
    \draw[ultra thick, blue!80!black] (2, 0) -- (4, 0);
    \draw[ultra thick, blue!80!black, dashed] (4, 0) -- (4, 3);
    \draw[ultra thick, blue!80!black] (4, 3) -- (5.5, 3);
    
    \draw[thick] (2, 0.1) -- (2, -0.1) node[below=2pt] {$2$};
    \draw[thick] (4, 0.1) -- (4, -0.1) node[below=2pt] {$4$};
    \draw[thick] (0.1, 3) -- (-0.1, 3) node[left=2pt] {$3$};
    
    \node[align=left, anchor=west, draw, inner sep=8pt, fill=white] at (6.5, 2) {Square Wave\\[4pt] $T=4$\\[4pt] Amplitude $= 3$};
\end{tikzpicture}
\end{center}

\vspace{1em}
\textbf{Problem 2:} Determine $\mathcal{L}\{f\}$, where $f(t) = 3t$ for $0 < t < 3$, and $f(t)$ has period $T = 3$.

\textbf{Solution:}
The period is $T = 3$. The Laplace transform is:
\[ \mathcal{L}\{f\}(s) = \frac{1}{1 - e^{-3s}} \int_0^3 3t e^{-st} \,dt \]
We evaluate the integral using integration by parts. Let $u = 3t$ and $dv = e^{-st} \,dt$. Then $du = 3 \,dt$ and $v = -\frac{1}{s}e^{-st}$.
\begin{align*}
\int_0^3 3t e^{-st} \,dt &= \left[ -\frac{3t}{s} e^{-st} \right]_0^3 - \int_0^3 -\frac{3}{s} e^{-st} \,dt \\
&= \left( -\frac{9}{s} e^{-3s} - 0 \right) - \left[ \frac{3}{s^2} e^{-st} \right]_0^3 \\
&= -\frac{9}{s} e^{-3s} - \left( \frac{3}{s^2} e^{-3s} - \frac{3}{s^2} \right) \\
&= \frac{3}{s^2} - \frac{9}{s} e^{-3s} - \frac{3}{s^2} e^{-3s}
\end{align*}
Substituting this back into the periodic transform formula:
\[ \mathcal{L}\{f\}(s) = \frac{3 - 9s e^{-3s} - 3 e^{-3s}}{s^2(1 - e^{-3s})} \]

\vspace{1em}
\textbf{Problem 3:} Determine $\mathcal{L}\{f\}$, where $f(t)$ is a triangular wave defined by $f(t) = 2t$ for $0 < t < 1$, and $f(t) = 4 - 2t$ for $1 < t < 2$, with period $T = 2$.

\textbf{Solution:}
The period is $T = 2$. We split the integral over the two sub-intervals:
\[ \mathcal{L}\{f\}(s) = \frac{1}{1 - e^{-2s}} \left( \int_0^1 2t e^{-st} \,dt + \int_1^2 (4 - 2t) e^{-st} \,dt \right) \]
First integral (using integration by parts):
\[ \int_0^1 2t e^{-st} \,dt = \left[ -\frac{2t}{s} e^{-st} - \frac{2}{s^2} e^{-st} \right]_0^1 = -\frac{2}{s} e^{-s} - \frac{2}{s^2} e^{-s} + \frac{2}{s^2} \]
Second integral. Let $u = t - 1$, so $t = u + 1$ and $dt = du$. The limits change from $0$ to $1$.
\begin{align*}
\int_1^2 (4 - 2t) e^{-st} \,dt &= \int_0^1 (4 - 2(u+1)) e^{-s(u+1)} \,du \\
&= e^{-s} \int_0^1 (2 - 2u) e^{-su} \,du \\
&= e^{-s} \left( \int_0^1 2 e^{-su} \,du - \int_0^1 2u e^{-su} \,du \right) \\
&= e^{-s} \left( \left[ -\frac{2}{s} e^{-su} \right]_0^1 - \left[ -\frac{2u}{s} e^{-su} - \frac{2}{s^2} e^{-su} \right]_0^1 \right) \\
&= e^{-s} \left( -\frac{2}{s} e^{-s} + \frac{2}{s} - \left( -\frac{2}{s} e^{-s} - \frac{2}{s^2} e^{-s} + \frac{2}{s^2} \right) \right) \\
&= e^{-s} \left( \frac{2}{s} + \frac{2}{s^2} e^{-s} - \frac{2}{s^2} \right) = \frac{2}{s} e^{-s} + \frac{2}{s^2} e^{-2s} - \frac{2}{s^2} e^{-s}
\end{align*}
Adding the two integrals together:
\[ \left( -\frac{2}{s} e^{-s} - \frac{2}{s^2} e^{-s} + \frac{2}{s^2} \right) + \left( \frac{2}{s} e^{-s} + \frac{2}{s^2} e^{-2s} - \frac{2}{s^2} e^{-s} \right) = \frac{2}{s^2} - \frac{4}{s^2} e^{-s} + \frac{2}{s^2} e^{-2s} = \frac{2(1 - e^{-s})^2}{s^2} \]
Divide by $1 - e^{-2s} = (1 - e^{-s})(1 + e^{-s})$:
\[ \mathcal{L}\{f\}(s) = \frac{2(1 - e^{-s})^2}{s^2(1 - e^{-s})(1 + e^{-s})} = \frac{2(1 - e^{-s})}{s^2(1 + e^{-s})} = \frac{2}{s^2} \tanh\left(\frac{s}{2}\right) \]

\vspace{1em}
\textbf{Problem 4:} Determine $\mathcal{L}\{f\}$, where $f(t)$ is a half-rectified sine wave defined by $f(t) = \sin(3t)$ for $0 < t < \pi/3$, and $f(t) = 0$ for $\pi/3 < t < 2\pi/3$, with period $T = 2\pi/3$.

\textbf{Solution:}
The period is $T = 2\pi/3$.
\[ \mathcal{L}\{f\}(s) = \frac{1}{1 - e^{-2\pi s / 3}} \int_0^{\pi/3} e^{-st} \sin(3t) \,dt \]
Using the standard integration formula $\int e^{at} \sin(bt) \,dt = \frac{e^{at}}{a^2+b^2} (a \sin(bt) - b \cos(bt))$ with $a = -s$ and $b = 3$:
\begin{align*}
\int_0^{\pi/3} e^{-st} \sin(3t) \,dt &= \left[ \frac{e^{-st}}{s^2 + 9} (-s \sin(3t) - 3 \cos(3t)) \right]_0^{\pi/3} \\
&= \frac{e^{-\pi s / 3}}{s^2 + 9} (0 - 3(-1)) - \frac{1}{s^2 + 9} (0 - 3(1)) \\
&= \frac{3 e^{-\pi s / 3} + 3}{s^2 + 9} = \frac{3(1 + e^{-\pi s / 3})}{s^2 + 9}
\end{align*}
Substitute into the transform formula:
\[ \mathcal{L}\{f\}(s) = \frac{3(1 + e^{-\pi s / 3})}{(s^2 + 9)(1 - e^{-2\pi s / 3})} \]
Since $1 - e^{-2\pi s / 3} = (1 - e^{-\pi s / 3})(1 + e^{-\pi s / 3})$, we cancel the common factor:
\[ \mathcal{L}\{f\}(s) = \frac{3}{(s^2 + 9)(1 - e^{-\pi s / 3})} \]

\vspace{1em}
\textbf{Problem 5:} Determine $\mathcal{L}\{f\}$, where $f(t) = |\cos(t)|$.

\textbf{Solution:}
The function $f(t) = |\cos(t)|$ is a full-rectified cosine wave. The fundamental period of $|\cos(t)|$ is $T = \pi$.
\[ \mathcal{L}\{f\}(s) = \frac{1}{1 - e^{-\pi s}} \int_0^\pi e^{-st} |\cos(t)| \,dt \]
Since $\cos(t)$ is positive on $(0, \pi/2)$ and negative on $(\pi/2, \pi)$, we split the integral:
\[ \int_0^\pi e^{-st} |\cos(t)| \,dt = \int_0^{\pi/2} e^{-st} \cos(t) \,dt - \int_{\pi/2}^\pi e^{-st} \cos(t) \,dt \]
Using $\int e^{-st} \cos(t) \,dt = \frac{e^{-st}}{s^2+1} (-s \cos(t) + \sin(t))$:
First integral:
\[ \left[ \frac{e^{-st}}{s^2+1} (-s \cos(t) + \sin(t)) \right]_0^{\pi/2} = \frac{e^{-\pi s / 2}}{s^2+1}(1) - \frac{1}{s^2+1}(-s) = \frac{e^{-\pi s / 2} + s}{s^2+1} \]
Second integral:
\[ \left[ \frac{e^{-st}}{s^2+1} (-s \cos(t) + \sin(t)) \right]_{\pi/2}^\pi = \frac{e^{-\pi s}}{s^2+1}(s) - \frac{e^{-\pi s / 2}}{s^2+1}(1) = \frac{s e^{-\pi s} - e^{-\pi s / 2}}{s^2+1} \]
Subtracting the second from the first:
\[ \frac{e^{-\pi s / 2} + s - s e^{-\pi s} + e^{-\pi s / 2}}{s^2+1} = \frac{s(1 - e^{-\pi s}) + 2e^{-\pi s / 2}}{s^2+1} \]
Thus, the Laplace transform is:
\[ \mathcal{L}\{f\}(s) = \frac{s(1 - e^{-\pi s}) + 2e^{-\pi s / 2}}{(s^2+1)(1 - e^{-\pi s})} = \frac{s}{s^2+1} + \frac{2e^{-\pi s / 2}}{(s^2+1)(1 - e^{-\pi s})} \]

\vspace{1em}
\textbf{Problem 6:} Determine $\mathcal{L}\{f\}$, where $f(t) = e^{2t}$ for $0 < t < 1$, with period $T = 1$.

\textbf{Solution:}
The period is $T = 1$.
\[ \mathcal{L}\{f\}(s) = \frac{1}{1 - e^{-s}} \int_0^1 e^{-st} e^{2t} \,dt = \frac{1}{1 - e^{-s}} \int_0^1 e^{(2-s)t} \,dt \]
Evaluating the integral:
\[ \int_0^1 e^{(2-s)t} \,dt = \left[ \frac{1}{2-s} e^{(2-s)t} \right]_0^1 = \frac{e^{2-s} - 1}{2-s} \]
Thus, the Laplace transform is:
\[ \mathcal{L}\{f\}(s) = \frac{e^{2-s} - 1}{(2-s)(1 - e^{-s})} \]

\vspace{1em}
\textbf{Problem 7:} Determine $\mathcal{L}\{f\}$, where $f(t) = t^2$ for $0 < t < 2$, with period $T = 2$.

\textbf{Solution:}
The period is $T = 2$.
\[ \mathcal{L}\{f\}(s) = \frac{1}{1 - e^{-2s}} \int_0^2 t^2 e^{-st} \,dt \]
Using tabular integration by parts for $\int t^2 e^{-st} \,dt$:
\begin{align*}
\int_0^2 t^2 e^{-st} \,dt &= \left[ t^2 \left(-\frac{1}{s}e^{-st}\right) - 2t \left(\frac{1}{s^2}e^{-st}\right) + 2 \left(-\frac{1}{s^3}e^{-st}\right) \right]_0^2 \\
&= \left( -\frac{4}{s}e^{-2s} - \frac{4}{s^2}e^{-2s} - \frac{2}{s^3}e^{-2s} \right) - \left( 0 - 0 - \frac{2}{s^3} \right) \\
&= \frac{2}{s^3} - e^{-2s} \left( \frac{4}{s} + \frac{4}{s^2} + \frac{2}{s^3} \right)
\end{align*}
Thus, the transform is:
\[ \mathcal{L}\{f\}(s) = \frac{\frac{2}{s^3} - e^{-2s} \left( \frac{4}{s} + \frac{4}{s^2} + \frac{2}{s^3} \right)}{1 - e^{-2s}} \]

\vspace{1em}
\textbf{Problem 8:} Determine $\mathcal{L}\{f\}$, where $f(t) = 1$ for $0 < t < a$, and $f(t) = -1$ for $a < t < 2a$, with period $T = 2a$.

\textbf{Solution:}
The period is $T = 2a$.
\[ \mathcal{L}\{f\}(s) = \frac{1}{1 - e^{-2as}} \left( \int_0^a 1 \cdot e^{-st} \,dt + \int_a^{2a} (-1) \cdot e^{-st} \,dt \right) \]
Evaluate the integrals:
\[ \int_0^a e^{-st} \,dt = \frac{1 - e^{-as}}{s} \]
\[ \int_a^{2a} -e^{-st} \,dt = \left[ \frac{1}{s} e^{-st} \right]_a^{2a} = \frac{e^{-2as} - e^{-as}}{s} \]
Adding them together:
\[ \frac{1 - e^{-as} + e^{-2as} - e^{-as}}{s} = \frac{1 - 2e^{-as} + e^{-2as}}{s} = \frac{(1 - e^{-as})^2}{s} \]
Substitute into the formula:
\[ \mathcal{L}\{f\}(s) = \frac{(1 - e^{-as})^2}{s(1 - e^{-2as})} = \frac{(1 - e^{-as})^2}{s(1 - e^{-as})(1 + e^{-as})} = \frac{1 - e^{-as}}{s(1 + e^{-as})} = \frac{1}{s} \tanh\left(\frac{as}{2}\right) \]

\vspace{1em}
\textbf{Problem 9:} Determine $\mathcal{L}\{f\}$, where $f(t) = |\sin(t)|$ (Full-rectified sine wave).

\textbf{Solution:}
The fundamental period of $|\sin(t)|$ is $T = \pi$.
\[ \mathcal{L}\{f\}(s) = \frac{1}{1 - e^{-\pi s}} \int_0^\pi e^{-st} \sin(t) \,dt \]
Using the standard integration formula:
\begin{align*}
\int_0^\pi e^{-st} \sin(t) \,dt &= \left[ \frac{e^{-st}}{s^2 + 1} (-s \sin(t) - \cos(t)) \right]_0^\pi \\
&= \frac{e^{-\pi s}}{s^2 + 1} (0 - (-1)) - \frac{1}{s^2 + 1} (0 - 1) \\
&= \frac{e^{-\pi s} + 1}{s^2 + 1}
\end{align*}
Substitute into the transform formula:
\[ \mathcal{L}\{f\}(s) = \frac{1 + e^{-\pi s}}{(s^2 + 1)(1 - e^{-\pi s})} = \frac{1}{s^2 + 1} \coth\left(\frac{\pi s}{2}\right) \]

\vspace{1em}
\textbf{Problem 10:} Determine $\mathcal{L}\{f\}$, where $f(t) = \sinh(t)$ for $0 < t < 2$, with period $T = 2$.

\textbf{Solution:}
The period is $T = 2$. We use the definition $\sinh(t) = \frac{e^t - e^{-t}}{2}$.
\[ \mathcal{L}\{f\}(s) = \frac{1}{1 - e^{-2s}} \int_0^2 e^{-st} \left( \frac{e^t - e^{-t}}{2} \right) \,dt = \frac{1}{2(1 - e^{-2s})} \int_0^2 \left( e^{(1-s)t} - e^{-(1+s)t} \right) \,dt \]
Evaluating the integral:
\begin{align*}
\int_0^2 \left( e^{(1-s)t} - e^{-(1+s)t} \right) \,dt &= \left[ \frac{e^{(1-s)t}}{1-s} + \frac{e^{-(1+s)t}}{1+s} \right]_0^2 \\
&= \left( \frac{e^{2(1-s)}}{1-s} + \frac{e^{-2(1+s)}}{1+s} \right) - \left( \frac{1}{1-s} + \frac{1}{1+s} \right) \\
&= \frac{e^{2-2s} - 1}{1-s} + \frac{e^{-2-2s} - 1}{1+s}
\end{align*}
Thus, the transform is:
\[ \mathcal{L}\{f\}(s) = \frac{1}{2(1 - e^{-2s})} \left( \frac{e^{2-2s} - 1}{1-s} + \frac{e^{-2-2s} - 1}{1+s} \right) \]

\vspace{1em}
\textbf{Problem 11:} Determine $\mathcal{L}\{f\}$, where $f(t) = t e^{-t}$ for $0 < t < 1$, with period $T = 1$.

\textbf{Solution:}
The period is $T = 1$.
\[ \mathcal{L}\{f\}(s) = \frac{1}{1 - e^{-s}} \int_0^1 t e^{-t} e^{-st} \,dt = \frac{1}{1 - e^{-s}} \int_0^1 t e^{-(s+1)t} \,dt \]
Using integration by parts with $u = t$ and $dv = e^{-(s+1)t} \,dt$:
\begin{align*}
\int_0^1 t e^{-(s+1)t} \,dt &= \left[ -\frac{t}{s+1} e^{-(s+1)t} - \frac{1}{(s+1)^2} e^{-(s+1)t} \right]_0^1 \\
&= \left( -\frac{1}{s+1} e^{-(s+1)} - \frac{1}{(s+1)^2} e^{-(s+1)} \right) - \left( 0 - \frac{1}{(s+1)^2} \right) \\
&= \frac{1}{(s+1)^2} - e^{-(s+1)} \left( \frac{1}{s+1} + \frac{1}{(s+1)^2} \right)
\end{align*}
Thus, the transform is:
\[ \mathcal{L}\{f\}(s) = \frac{\frac{1}{(s+1)^2} - e^{-(s+1)} \left( \frac{s+2}{(s+1)^2} \right)}{1 - e^{-s}} \]

\vspace{1em}
\textbf{Problem 12:} Determine $\mathcal{L}\{f\}$, where $f(t) = 1 - e^{-t}$ for $0 < t < 2$, with period $T = 2$.

\textbf{Solution:}
The period is $T = 2$.
\[ \mathcal{L}\{f\}(s) = \frac{1}{1 - e^{-2s}} \int_0^2 (1 - e^{-t}) e^{-st} \,dt = \frac{1}{1 - e^{-2s}} \int_0^2 \left( e^{-st} - e^{-(s+1)t} \right) \,dt \]
Evaluating the integral:
\begin{align*}
\int_0^2 \left( e^{-st} - e^{-(s+1)t} \right) \,dt &= \left[ -\frac{1}{s} e^{-st} + \frac{1}{s+1} e^{-(s+1)t} \right]_0^2 \\
&= \left( -\frac{1}{s} e^{-2s} + \frac{1}{s+1} e^{-2(s+1)} \right) - \left( -\frac{1}{s} + \frac{1}{s+1} \right) \\
&= \frac{1 - e^{-2s}}{s} - \frac{1 - e^{-2(s+1)}}{s+1}
\end{align*}
Thus, the transform is:
\[ \mathcal{L}\{f\}(s) = \frac{1}{s} - \frac{1 - e^{-2(s+1)}}{(s+1)(1 - e^{-2s})} \]

\vspace{1em}
\textbf{Problem 13:} Find the inverse Laplace transform of $F(s) = \frac{1}{s(1 - e^{-4s})}$. Express the result as an infinite series of step functions.

\textbf{Solution:}
We expand the denominator as a geometric series:
\[ \frac{1}{1 - e^{-4s}} = \sum_{n=0}^\infty e^{-4ns} \]
Multiply by $\frac{1}{s}$:
\[ F(s) = \sum_{n=0}^\infty \frac{e^{-4ns}}{s} \]
We know that $\mathcal{L}^{-1} \left\{ \frac{1}{s} \right\} = 1$. By the Second Translation Theorem, $\mathcal{L}^{-1}\{e^{-as}G(s)\} = g(t - a)u(t - a)$. Applying this to each term with $a = 4n$:
\[ \mathcal{L}^{-1} \left\{ \frac{e^{-4ns}}{s} \right\} = 1 \cdot u(t - 4n) \]
Summing these inverse transforms yields a staircase function:
\[ f(t) = \sum_{n=0}^\infty u(t - 4n) \]

\vspace{1em}
\textbf{Problem 14:} Find the inverse Laplace transform of $F(s) = \frac{1}{(s^2 + 9)(1 - e^{-\pi s})}$.

\textbf{Solution:}
Expand the denominator as a geometric series:
\[ F(s) = \frac{1}{s^2 + 9} \sum_{n=0}^\infty e^{-n\pi s} = \frac{1}{3} \sum_{n=0}^\infty \frac{3 e^{-n\pi s}}{s^2 + 9} \]
We know that $\mathcal{L}^{-1} \left\{ \frac{3}{s^2 + 9} \right\} = \sin(3t)$. Applying the Second Translation Theorem:
\[ \mathcal{L}^{-1} \left\{ \frac{3 e^{-n\pi s}}{s^2 + 9} \right\} = \sin(3(t - n\pi)) u(t - n\pi) \]
Using the trigonometric identity $\sin(3t - 3n\pi) = \sin(3t)\cos(3n\pi) - \cos(3t)\sin(3n\pi) = \sin(3t)(-1)^{3n} = (-1)^n \sin(3t)$:
\[ f(t) = \frac{1}{3} \sum_{n=0}^\infty (-1)^n \sin(3t) u(t - n\pi) \]

\vspace{1em}
\textbf{Problem 15:} Find the inverse Laplace transform of $F(s) = \frac{e^{-2s}}{(s + 3)(1 - e^{-2s})}$.

\textbf{Solution:}
Expand the denominator as a geometric series:
\[ F(s) = \frac{e^{-2s}}{s + 3} \sum_{n=0}^\infty e^{-2ns} = \sum_{n=0}^\infty \frac{e^{-2(n+1)s}}{s + 3} \]
Let $k = n + 1$. The sum becomes $\sum_{k=1}^\infty \frac{e^{-2ks}}{s + 3}$.
We know that $\mathcal{L}^{-1} \left\{ \frac{1}{s + 3} \right\} = e^{-3t}$. Applying the Second Translation Theorem:
\[ f(t) = \sum_{k=1}^\infty e^{-3(t - 2k)} u(t - 2k) \]

\vspace{1em}
\textbf{Problem 16:} Find the inverse Laplace transform of $F(s) = \frac{1 - e^{-s}}{s^2(1 - e^{-2s})}$.

\textbf{Solution:}
Notice that $1 - e^{-2s} = (1 - e^{-s})(1 + e^{-s})$. Canceling the common term $1 - e^{-s}$:
\[ F(s) = \frac{1}{s^2(1 + e^{-s})} \]
Expand $\frac{1}{1 + e^{-s}}$ as a geometric series with ratio $-e^{-s}$:
\[ F(s) = \frac{1}{s^2} \sum_{n=0}^\infty (-1)^n e^{-ns} = \sum_{n=0}^\infty (-1)^n \frac{e^{-ns}}{s^2} \]
We know that $\mathcal{L}^{-1} \left\{ \frac{1}{s^2} \right\} = t$. Applying the Second Translation Theorem:
\[ f(t) = \sum_{n=0}^\infty (-1)^n (t - n) u(t - n) \]
This represents a triangular wave.

\vspace{1em}
\textbf{Problem 17:} Show that $\mathcal{L}^{-1} \left\{ \frac{1}{s(1 + e^{-s})} \right\}$ represents an alternating square wave.

\textbf{Solution:}
Expand the denominator as a geometric series:
\[ F(s) = \frac{1}{s} \sum_{n=0}^\infty (-1)^n e^{-ns} = \sum_{n=0}^\infty (-1)^n \frac{e^{-ns}}{s} \]
Taking the inverse transform term-by-term:
\[ f(t) = \sum_{n=0}^\infty (-1)^n u(t - n) \]
Let's evaluate this sum on intervals of length 1:
- For $0 < t < 1$, only the $n = 0$ term is active: $f(t) = 1$.
- For $1 < t < 2$, the $n = 0$ and $n = 1$ terms are active: $f(t) = 1 - 1 = 0$.
- For $2 < t < 3$, the $n = 0, 1, 2$ terms are active: $f(t) = 1 - 1 + 1 = 1$.
This perfectly matches the definition of a square wave that alternates between 1 and 0 with period 2.

\vspace{1em}
\textbf{Problem 18:} Evaluate $\Gamma(9/2)$ using the recursive property of the Gamma function.

\textbf{Solution:}
The recursive property is $\Gamma(r + 1) = r\Gamma(r)$. We also know $\Gamma(1/2) = \sqrt{\pi}$.
\begin{align*}
\Gamma(9/2) &= \frac{7}{2}\Gamma(7/2) \\
&= \frac{7}{2} \cdot \frac{5}{2}\Gamma(5/2) \\
&= \frac{7}{2} \cdot \frac{5}{2} \cdot \frac{3}{2}\Gamma(3/2) \\
&= \frac{7}{2} \cdot \frac{5}{2} \cdot \frac{3}{2} \cdot \frac{1}{2}\Gamma(1/2) \\
&= \frac{105}{16}\sqrt{\pi}
\end{align*}

\vspace{1em}
\textbf{Problem 19:} Evaluate the integral $\int_0^\infty x^4 e^{-x^2} \,dx$ by relating it to the Gamma function.

\textbf{Solution:}
Let $u = x^2$. Then $x = u^{1/2}$ and $dx = \frac{1}{2}u^{-1/2} \,du$. As $x \to 0$, $u \to 0$, and as $x \to \infty$, $u \to \infty$. Substitute these into the integral:
\begin{align*}
\int_0^\infty x^4 e^{-x^2} \,dx &= \int_0^\infty (u^{1/2})^4 e^{-u} \left( \frac{1}{2}u^{-1/2} \right) \,du \\
&= \frac{1}{2} \int_0^\infty u^2 e^{-u} u^{-1/2} \,du \\
&= \frac{1}{2} \int_0^\infty e^{-u} u^{3/2} \,du
\end{align*}
By definition, $\Gamma(r) = \int_0^\infty e^{-u} u^{r-1} \,du$. Here, $r - 1 = 3/2$, so $r = 5/2$.
\[ \frac{1}{2} \Gamma(5/2) = \frac{1}{2} \left( \frac{3}{2} \cdot \frac{1}{2} \Gamma(1/2) \right) = \frac{1}{2} \left( \frac{3}{4} \sqrt{\pi} \right) = \frac{3}{8}\sqrt{\pi} \]

\vspace{1em}
\textbf{Problem 20:} Determine $\mathcal{L}\{t^{7/2}\}$.

\textbf{Solution:}
Using the formula $\mathcal{L}\{t^r\} = \frac{\Gamma(r+1)}{s^{r+1}}$ with $r = 7/2$:
\[ \mathcal{L}\{t^{7/2}\} = \frac{\Gamma(9/2)}{s^{9/2}} \]
From Problem 18, we know $\Gamma(9/2) = \frac{105}{16}\sqrt{\pi}$.
\[ \mathcal{L}\{t^{7/2}\} = \frac{105\sqrt{\pi}}{16s^{9/2}} \]

\vspace{1em}
\textbf{Problem 21:} Determine $\mathcal{L}\{t^{-1/2} e^{4t}\}$ using the First Translation Theorem.

\textbf{Solution:}
First, find the Laplace transform of $f(t) = t^{-1/2}$.
\[ \mathcal{L}\{t^{-1/2}\} = \frac{\Gamma(1/2)}{s^{1/2}} = \frac{\sqrt{\pi}}{\sqrt{s}} \]
By the First Translation Theorem, $\mathcal{L}\{e^{at} f(t)\} = F(s - a)$. Here $a = 4$:
\[ \mathcal{L}\{t^{-1/2} e^{4t}\} = \frac{\sqrt{\pi}}{\sqrt{s - 4}} \]

\vspace{1em}
\textbf{Problem 22:} Show that $\Gamma(n + 1/2) = \frac{(2n)! \sqrt{\pi}}{4^n n!}$ for any integer $n \ge 0$.

\textbf{Solution:}
We use the recursive relation $\Gamma(r + 1) = r\Gamma(r)$:
\begin{align*}
\Gamma(n + 1/2) &= (n - 1/2)(n - 3/2) \cdots (1/2)\Gamma(1/2) \\
&= \left(\frac{2n - 1}{2}\right) \left(\frac{2n - 3}{2}\right) \cdots \left(\frac{1}{2}\right) \sqrt{\pi} \\
&= \frac{1 \cdot 3 \cdot 5 \cdots (2n - 1)}{2^n} \sqrt{\pi}
\end{align*}
To express the numerator in terms of factorials, multiply the numerator and denominator by the product of even integers $2 \cdot 4 \cdot 6 \cdots (2n) = 2^n n!$:
\[ \frac{1 \cdot 3 \cdot 5 \cdots (2n - 1)}{2^n} \cdot \frac{2^n n!}{2^n n!} = \frac{(2n)!}{2^{2n} n!} = \frac{(2n)!}{4^n n!} \]
Thus, $\Gamma(n + 1/2) = \frac{(2n)! \sqrt{\pi}}{4^n n!}$.

\vspace{1em}
\textbf{Problem 23:} Find the inverse Laplace transform of $F(s) = s^{-9/2}$.

\textbf{Solution:}
We know that $\mathcal{L}\{t^r\} = \frac{\Gamma(r+1)}{s^{r+1}}$. Let $r + 1 = 9/2$, so $r = 7/2$.
\[ \mathcal{L}^{-1} \{s^{-9/2}\} = \frac{t^{7/2}}{\Gamma(9/2)} \]
From Problem 18, $\Gamma(9/2) = \frac{105}{16}\sqrt{\pi}$.
\[ \mathcal{L}^{-1} \{s^{-9/2}\} = \frac{16 t^{7/2}}{105\sqrt{\pi}} \]

\vspace{1em}
\textbf{Problem 24:} Find an expansion for $\mathcal{L} \left\{ \frac{\sinh(t)}{t} \right\}$ in powers of $1/s$ by transforming the Maclaurin series term-by-term.

\textbf{Solution:}
The Maclaurin series for $\sinh(t)$ is $\sum_{n=0}^\infty \frac{t^{2n+1}}{(2n+1)!}$.
Divide by $t$:
\[ f(t) = \frac{\sinh(t)}{t} = \sum_{n=0}^\infty \frac{t^{2n}}{(2n+1)!} \]
Transform term-by-term using $\mathcal{L}\{t^{2n}\} = \frac{(2n)!}{s^{2n+1}}$:
\[ \mathcal{L}\{f(t)\} = \sum_{n=0}^\infty \frac{1}{(2n+1)!} \frac{(2n)!}{s^{2n+1}} = \sum_{n=0}^\infty \frac{1}{(2n+1)s^{2n+1}} \]
This is the Maclaurin series for $\frac{1}{2} \ln\left(\frac{s+1}{s-1}\right)$.

\vspace{1em}
\textbf{Problem 25:} Find an expansion for $\mathcal{L} \{\cos(2\sqrt{t})\}$ by transforming its Maclaurin series term-by-term.

\textbf{Solution:}
The Maclaurin series for $\cos(x)$ is $\sum_{n=0}^\infty \frac{(-1)^n x^{2n}}{(2n)!}$. Substitute $x = 2\sqrt{t}$:
\[ \cos(2\sqrt{t}) = \sum_{n=0}^\infty \frac{(-1)^n (2\sqrt{t})^{2n}}{(2n)!} = \sum_{n=0}^\infty \frac{(-1)^n 4^n t^n}{(2n)!} \]
Transform term-by-term using $\mathcal{L}\{t^n\} = \frac{n!}{s^{n+1}}$:
\[ \mathcal{L}\{\cos(2\sqrt{t})\} = \sum_{n=0}^\infty \frac{(-1)^n 4^n n!}{(2n)! s^{n+1}} \]

\vspace{1em}
\textbf{Problem 26:} Find the inverse Laplace transform of $F(s) = \frac{1}{s} e^{-3/s}$ using a series expansion.

\textbf{Solution:}
The Maclaurin series for $e^x$ is $\sum_{n=0}^\infty \frac{x^n}{n!}$. Letting $x = -3/s$:
\[ e^{-3/s} = \sum_{n=0}^\infty \frac{(-1)^n 3^n}{n! s^n} \]
Multiply by $1/s$:
\[ F(s) = \frac{1}{s} e^{-3/s} = \sum_{n=0}^\infty \frac{(-1)^n 3^n}{n! s^{n+1}} \]
Take the inverse Laplace transform term-by-term using $\mathcal{L}^{-1} \left\{ \frac{1}{s^{n+1}} \right\} = \frac{t^n}{n!}$:
\[ f(t) = \sum_{n=0}^\infty \frac{(-1)^n 3^n t^n}{n! n!} = \sum_{n=0}^\infty \frac{(-1)^n (3t)^n}{(n!)^2} \]
This is related to the Bessel function of the first kind of order zero, $J_0(2\sqrt{3t})$.

\vspace{1em}
\textbf{Problem 27:} Solve the initial value problem $y' + 2y = f(t)$ with $y(0) = 0$, where $f(t)$ is a square wave defined by $f(t) = 1$ for $0 < t < 1$ and $f(t) = 0$ for $1 < t < 2$, with period $T = 2$.

\textbf{Solution:}
Take the Laplace transform of both sides. Since $y(0) = 0$, $\mathcal{L}\{y'\} = sY(s)$.
\[ sY(s) + 2Y(s) = F(s) \implies Y(s) = \frac{F(s)}{s + 2} \]
From Problem 8 (with $a=1$, shifted up and scaled), or by direct integration, the transform of this specific square wave is:
\[ F(s) = \frac{1}{1 - e^{-2s}} \int_0^1 e^{-st} \,dt = \frac{1 - e^{-s}}{s(1 - e^{-2s})} = \frac{1}{s(1 + e^{-s})} \]
Substitute $F(s)$ into $Y(s)$:
\[ Y(s) = \frac{1}{s(s + 2)(1 + e^{-s})} \]
Use partial fractions on the transfer function:
\[ \frac{1}{s(s + 2)} = \frac{1/2}{s} - \frac{1/2}{s + 2} \]
So, $Y(s) = \frac{1}{2} \left( \frac{1}{s} - \frac{1}{s + 2} \right) \frac{1}{1 + e^{-s}}$.
Expand $\frac{1}{1 + e^{-s}} = \sum_{n=0}^\infty (-1)^n e^{-ns}$.
\[ Y(s) = \frac{1}{2} \sum_{n=0}^\infty (-1)^n e^{-ns} \left( \frac{1}{s} - \frac{1}{s + 2} \right) \]
Take the inverse transform term-by-term using the Second Translation Theorem:
\[ y(t) = \frac{1}{2} \sum_{n=0}^\infty (-1)^n \left( 1 - e^{-2(t - n)} \right) u(t - n) \]

\vspace{1em}
\textbf{Problem 28:} Solve the initial value problem $y'' + y = \text{III}_{2\pi}(t)$ with $y(0) = 0, y'(0) = 0$, where $\text{III}_{2\pi}(t) = \sum_{n=0}^\infty \delta(t - 2n\pi)$ is the Dirac comb.

\textbf{Solution:}
Take the Laplace transform of both sides.
\[ s^2 Y(s) + Y(s) = \mathcal{L}\{\text{III}_{2\pi}(t)\} \]
The transform of the Dirac comb is:
\[ \mathcal{L} \left\{ \sum_{n=0}^\infty \delta(t - 2n\pi) \right\} = \sum_{n=0}^\infty e^{-2n\pi s} = \frac{1}{1 - e^{-2\pi s}} \]
So, $Y(s) = \frac{1}{(s^2 + 1)(1 - e^{-2\pi s})}$.
From Problem 14 (with $\omega = 1$ and period $2\pi$), we expand the denominator:
\[ Y(s) = \frac{1}{s^2 + 1} \sum_{n=0}^\infty e^{-2n\pi s} \]
Take the inverse transform term-by-term:
\[ y(t) = \sum_{n=0}^\infty \sin(t - 2n\pi) u(t - 2n\pi) \]
Since $\sin(t - 2n\pi) = \sin(t)$, this simplifies to:
\[ y(t) = \sin(t) \sum_{n=0}^\infty u(t - 2n\pi) \]
This represents a system experiencing resonance due to periodic impulses perfectly timed with its natural frequency.

\vspace{1em}
\textbf{Problem 29:} Use the derivative shortcut $\mathcal{L}\{f'\} = sF(s) - f(0)$ to find the Laplace transform of the sawtooth wave $f(t) = t$ for $0 < t < a$, period $a$.

\textbf{Solution:}
The function is $f(t) = t - a \sum_{n=1}^\infty u(t - na)$. Its derivative is $f'(t) = 1 - a \sum_{n=1}^\infty \delta(t - na)$.
Take the Laplace transform of $f'(t)$:
\[ \mathcal{L}\{f'\} = \frac{1}{s} - a \sum_{n=1}^\infty e^{-nas} = \frac{1}{s} - \frac{a e^{-as}}{1 - e^{-as}} \]
Using the derivative property $\mathcal{L}\{f'\} = sF(s) - f(0)$, and noting $f(0) = 0$:
\[ sF(s) = \frac{1}{s} - \frac{a e^{-as}}{1 - e^{-as}} \]
Divide by $s$:
\[ F(s) = \frac{1}{s^2} - \frac{a e^{-as}}{s(1 - e^{-as})} = \frac{1 - e^{-as} - as e^{-as}}{s^2(1 - e^{-as})} \]
This perfectly matches the result obtained via integration by parts.

\vspace{1em}
\textbf{Problem 30:} \textbf{Flawed Problem:} Find the Laplace transform of $f(t) = t^{3/2}$. \\
\textit{Flawed Solution:} Using the formula $\mathcal{L}\{t^r\} = \frac{\Gamma(r+1)}{s^{r+1}}$, we have $r = 3/2$.
\[ \mathcal{L}\{t^{3/2}\} = \frac{\Gamma(3/2)}{s^{3/2}} \]
We know $\Gamma(3/2) = \frac{1}{2}\Gamma(1/2) = \frac{1}{2}\sqrt{\pi}$. Therefore, the final answer is $\frac{\sqrt{\pi}}{2s^{3/2}}$. \\
\textbf{Task:} Identify the error, explain why it is wrong in one sentence, and provide the correct final answer.

\textbf{Grader's Answer Key:}
\textbf{Correction:} The formula requires the denominator to be $s^{r+1}$, so for $r = 3/2$, the denominator must be $s^{5/2}$, not $s^{3/2}$. Furthermore, the numerator should be $\Gamma(5/2)$, not $\Gamma(3/2)$.
The correct setup is $\mathcal{L}\{t^{3/2}\} = \frac{\Gamma(5/2)}{s^{5/2}}$.
Evaluating $\Gamma(5/2) = \frac{3}{2}\Gamma(3/2) = \frac{3}{4}\sqrt{\pi}$.
\textbf{Final Answer:} $\mathcal{L}\{t^{3/2}\} = \frac{3\sqrt{\pi}}{4s^{5/2}}$.

\end{document}